%% ============================================================
%%  Capítulo 12 – Riesgos y Decisiones Sensibles Residuales
%% ============================================================
\chapter{Riesgos y Decisiones Sensibles Residuales}
\label{chap:riesgos}

\section{Estado Actualizado de los Riesgos Identificados en el Entregable~2}
\label{sec:risk_actualizacion}

\begin{riskbox}[R-01: Disponibilidad del SoC -- Estado: Mitigado]
  \textbf{Acción tomada en E3:} Se seleccionó el Amlogic A311D2 con el
  esquemático de referencia público de Khadas VIM4 como guía de diseño. La
  disponibilidad del A311D2 se verificó en distribuidores de referencia
  (Digi-Key, Mouser). Se identifica como alternativa el Rockchip RK3399.\\
  \textbf{Riesgo residual para E4:} Bajo. El A311D2 tiene disponibilidad
  en distribuidores estándar al momento del diseño.
\end{riskbox}

\begin{riskbox}[R-02: EMI entre Subsistemas -- Estado: Mitigado por Partición]
  \textbf{Acción tomada en E3:} La partición en dos PCBs (Capítulo~\ref{chap:particion})
  aisla físicamente los reguladores Buck de las rutas de audio y señales
  de alta velocidad. En la PCB-CTRL, las zonas analógica, digital y de potencia
  están definidas en el esquemático con criterios de SI documentados
  (Capítulo~\ref{chap:si}).\\
  \textbf{Riesgo residual para E4:} Medio. La efectividad del aislamiento depende
  de la implementación correcta de las zonas en el layout. Las restricciones de SI
  del Capítulo~\ref{chap:si} deben trasladarse fielmente a las reglas de ruteo.
\end{riskbox}

\begin{riskbox}[R-03: Disipación Térmica -- Estado: Pendiente de Verificación]
  \textbf{Acción tomada en E3:} Se definió la estrategia de disipación pasiva
  mediante planos de cobre y vías térmicas bajo el SoC. La estimación del E2
  (19\,W totales) sigue vigente con los componentes seleccionados.\\
  \textbf{Riesgo residual para E4:} Medio-alto. Se requiere confirmación mediante
  simulación térmica (CFD básico) o cálculo analítico con el Theta\_JB del A311D2
  antes de cerrar el layout (ver Anexo~\ref{chap:anexo_si}).
\end{riskbox}

\begin{riskbox}[R-04: Latencia de Voz Fuera de Especificación -- Estado: No Aplica en E3]
  La latencia de voz depende principalmente del stack de software (OVOS,
  PipeWire, drivers de audio). El hardware está dimensionado con la potencia
  de cómputo del A311D2 (big.LITTLE) para cumplir con $<$2\,s de latencia
  de respuesta \reqref{REQ-F-02}. La validación final corresponde a la fase
  de integración de software (fuera del alcance del proyecto académico).
\end{riskbox}

\begin{riskbox}[R-05: Complejidad Normativa -- Estado: Documentado]
  Los requisitos de fabricante y SI están documentados en los Capítulos~\ref{chap:fabricante}
  y~\ref{chap:si}, y configurados como reglas DRC en Altium. Las distancias de
  clearance de IEC\,62368-1 \reqref{REQ-E-01} están establecidas para la entrada de
  12\,V. La verificación en prototipo físico corresponde al Entregable~4 y más allá.
\end{riskbox}

\section{Nuevos Riesgos Identificados en Este Entregable}
\label{sec:risk_nuevos}

\begin{riskbox}[R-06 (Crítico): Consistencia del Pinout del A311D2]
  \textbf{Descripción:} El componente del A311D2 en Altium Designer fue construido
  manualmente a partir del esquemático de referencia de Khadas, ya que no se dispone
  de un modelo de biblioteca oficial del fabricante. La configuración de todos los
  pines como tipo ``Pasivo'' suprime la validación automática de ERC, trasladando
  la responsabilidad de verificación a una revisión manual.\\
  \textbf{Probabilidad:} Alta (el error humano en la asignación de pines de un
  BGA-1056 es estadísticamente probable en al menos un subconjunto de pines).\\
  \textbf{Impacto:} Crítico. Un error de netlist en el SoC no detectado hasta el
  prototipo implicaría una revisión completa de la PCB-CTRL.\\
  \textbf{Mitigación obligatoria antes del layout:} Realizar una revisión cruzada
  sistemática del netlist generado por Altium contra el esquemático fuente de
  Khadas VIM4, verificando al menos: (a)~pines de alimentación y GND,
  (b)~interfaz MIPI DSI, (c)~buses I2S/PDM, (d)~buses I2C y SDIO.
\end{riskbox}

\begin{riskbox}[R-07: Dispersión de Pines de Voltaje en el BGA-1056]
  \textbf{Descripción:} Los pines de alimentación del A311D2 (VDD\_CPU-A,
  VDD\_CPU-B, VDD\_GPU, VDD\_NPU, etc.) están físicamente dispersos a lo largo
  del BGA. Esta dispersión complica la distribución de los planos de potencia
  en el layout y podría requerir una capa adicional exclusivamente para
  distribución de potencia.\\
  \textbf{Mitigación:} Durante el layout del E4, analizar la distribución de los
  pines de potencia del BGA y determinar si el stack-up de 6~capas es suficiente
  o si se requieren 8~capas. Decidir antes de iniciar el fan-out.
\end{riskbox}

\begin{riskbox}[R-08: Componentes con Riesgo de EOL o Disponibilidad Restringida]
  \textbf{Descripción:} Algunos componentes del BOM presentan disponibilidad
  limitada o estado End-of-Life (EOL) en los distribuidores consultados.
  Esto afecta principalmente a ciertos componentes pasivos de valores específicos
  para las redes de feedback de los reguladores.\\
  \textbf{Mitigación:} Identificar alternativas pin-compatible para los componentes
  en riesgo. Verificar stock disponible en Digi-Key y Mouser antes del cierre
  definitivo del BOM.
\end{riskbox}

\section{Decisiones Técnicas Sensibles para el Enrutamiento (E4)}
\label{sec:risk_e4}

Los siguientes puntos deben resolverse al inicio del Entregable~4, antes de
comenzar el ruteo de cualquier señal crítica:

\begin{enumerate}
  \item \textbf{Revisión cruzada del netlist del A311D2} (R-06): validación
    manual obligatoria contra el esquemático de Khadas.
  \item \textbf{Definición del fan-out del BGA-1056}: estrategia de escape de
    señales del SoC hacia las capas internas. Decidir uso de microvías (blind vias)
    si el proceso HDI de JLCPCB lo permite a costo aceptable.
  \item \textbf{Asignación de capas}: qué señales van en qué capa; confirmar
    que el stack-up de 6~capas es suficiente (R-07).
  \item \textbf{Ubicación de micrófonos MEMS}: posición física en la PCB-CTRL
    maximizando la distancia al amplificador y respetando las restricciones
    acústicas del enclosure.
  \item \textbf{Posición de la antena WiFi}: zona de exclusión de cobre bajo
    la antena del módulo WiFi/BT para no degradar la eficiencia de radiación.
\end{enumerate}
